\section{Diseño e implementación del software}
\subsection{Arquitectura modular}
\begin{figure}[H]\centering
\resizebox{\linewidth}{!}{\begin{tikzpicture}[>=latex,font=\small,
box/.style={draw=bluea,rounded corners=3mm,fill=bluea!7,minimum width=3.8cm,minimum height=1.7cm,align=center},
dat/.style={draw=teal,rounded corners=2mm,fill=teal!8,minimum width=4.2cm,minimum height=1.5cm,align=center},
arrow/.style={-{Latex[length=3mm]},thick,draw=navy}]

\node[box] (web) at (0,0) {\textbf{Navegador del votante}\\ HTML / JavaScript\\ Web Crypto API};
\node[box] (flask) at (6.0,0) {\textbf{Aplicación Flask}\\ sesión / Argon2id / CSRF\\ API de voto};
\node[dat] (voters) at (12.4,1.7) {\textbf{SQLite: voters}\\ usuario / hash Argon2id\\ marca de participación};
\node[dat] (ballots) at (12.4,-1.7) {\textbf{SQLite: ballots}\\ sobres cifrados / huellas\\ cadena SHA-256};
\node[box] (authority) at (18.9,-1.7) {\textbf{Autoridad honesta}\\ CLI: cierre / escrutinio\\ RSA privada / Ed25519};


\draw[arrow] (web) -- node[above,align=center]{HTTP local + CSRF\\sobre cifrado} (flask);
\draw[arrow] (flask) -- (voters);
\draw[arrow] (flask) -- (ballots);
\draw[arrow] (ballots) -- node[above]{lectura offline} (authority);
\draw[arrow,dashed] (authority.south) -- ++(0,-1.2) -| node[below,pos=.25]{resultados firmados y auditoría} (flask.south);

\node[draw=grayb,fill=gray!7,rounded corners=2mm,align=center,font=\footnotesize] at (10.2,-4.4) {Clave \textbf{pública} RSA en el navegador; clave \textbf{privada} solamente en el escrutinio local.\\ Las dos tablas SQLite comparten el mismo archivo físico: no existe anonimato fuerte.};

\end{tikzpicture}
}
\caption{Arquitectura lógica. El almacén SQLite contiene dos tablas separadas, no dos servidores independientes.}
\end{figure}

\subsection{Lenguajes y tecnologías}
\begin{longtable}{p{3.0cm}p{10.8cm}}
\toprule Tecnología & Justificación técnica\\\midrule\endhead
Python 3 & Facilita pruebas, modularidad e integración con bibliotecas de criptografía revisadas.\\
Flask & Rutas de aplicación web sencillas, sesiones firmadas y plantillas Jinja2.\\
JavaScript + Web Crypto & Cifra dentro del navegador; la selección no viaja en texto claro dentro del cuerpo de la petición. Web Crypto requiere un contexto seguro (HTTPS o localhost) \cite{mdn}.\\
SQLite & Almacenamiento transaccional de uso único y control de concurrencia mediante \code{BEGIN IMMEDIATE}.\\
\code{cryptography} & Implementaciones RSA-OAEP, AES-GCM y Ed25519 sin escribir primitivas manualmente.\\
\code{argon2-cffi} & Implementación Argon2id para almacenar hashes de contraseñas.\\
Pytest & Automatiza comprobación de casos normales, inválidos y concurrentes.\\
LaTeX y TikZ & Documentación reproducible con fórmulas y diagramas vectoriales editables.\\
\bottomrule
\end{longtable}

\subsection{Estructura de archivos}
\begin{center}\small
\begin{tabularx}{\linewidth}{p{5cm}X}
\toprule Archivo & Función\\\midrule
\code{app.py} & Rutas HTTP, autorización por sesión y protección CSRF.\\
\code{launcher.py} & Arranque web local en \code{127.0.0.1:5000}.\\
\code{manage.py} & Inicializar, cerrar, escrutar y verificar desde terminal.\\
\code{src/votoseguro/config.py} & Identificador electoral, candidatos y constantes.\\
\code{src/votoseguro/crypto.py} & Claves RSA, Ed25519, AES-GCM y funciones SHA-256.\\
\code{src/votoseguro/storage.py} & SQLite, voto único y cadena hash.\\
\code{src/votoseguro/services.py} & Argon2id, autenticación, auditoría y escrutinio.\\
\code{static/vote.js} & Cifrado de boleta con Web Crypto API en el cliente.\\
\code{templates/*.html} & Acceso, elección, comprobante, auditoría y resultados.\\
\code{tests/test\_election.py} & Pruebas de integridad y comportamiento de los protocolos.\\
\bottomrule
\end{tabularx}
\end{center}

\subsection{Base de datos y separación lógica}
La tabla \code{voters(username, password\_hash, has\_voted)} almacena credenciales de prueba y participación. La tabla \code{ballots} almacena sobres cifrados, nonces y hashes de integridad sin incluir el identificador del votante. \code{settings} conserva el estado abierto/cerrado.

\textbf{Limitación:} aunque no existe una columna explícita que relacione votante y selección, el mismo sistema web sigue observando autenticaciones, solicitudes y tiempos. Sin tecnologías de credenciales ciegas, redes anónimas y controles adicionales, no se garantiza desvinculación criptográfica fuerte.

\subsection{Pseudocódigo del registro de voto}
\begin{lstlisting}
FUNCION registrar_voto(usuario, sobre):
    validar_estructura(sobre)
    comenzar_transaccion_inmediata()
    SI eleccion_cerrada(): RECHAZAR
    SI actualizar_votante_si_no_voto(usuario) != 1:
        RECHAZAR
    recibo = SHA256(JSON_canonico(sobre))
    anterior = ultimo_hash_o_genesis()
    cadena = SHA256(anterior || indice || recibo)
    guardar_sobre(sobre, recibo, anterior, cadena)
    CONFIRMAR_TRANSACCION
    RETORNAR recibo
FIN FUNCION
\end{lstlisting}
\subsection{Diagrama de flujo: emisión}
\begin{figure}[H]\centering
\resizebox{0.85\linewidth}{!}{\begin{tikzpicture}[>=latex,font=\small,node distance=0.6cm,
boxstep/.style={rectangle,draw=bluea,fill=bluea!7,rounded corners=2mm,text width=6cm,align=center,minimum height=.85cm},
decision/.style={diamond,draw=teal,fill=teal!8,aspect=2.6,text width=3.5cm,align=center,inner sep=1pt},
edge/.style={-{Latex[length=2.5mm]},thick,draw=navy}]

\node[boxstep] (start) {Inicio: solicitar boleta};
\node[decision,below=of start] (login) {¿Sesión y CSRF válidos?};
\node[boxstep,below=0.95cm of login] (enc) {En el navegador: generar clave AES y nonce; cifrar con AES-GCM y RSA-OAEP};
\node[boxstep,below=of enc] (db) {Enviar sobre cifrado e iniciar transacción SQLite inmediata};
\node[decision,below=0.7cm of db] (elig) {¿Cuenta sin voto previo y elección abierta?};
\node[boxstep,below=1.1cm of elig] (save) {Marcar participación, guardar sobre y calcular cadena SHA-256};
\node[boxstep,below=of save] (end) {Confirmar transacción y devolver comprobante};
\node[boxstep,right=2.7cm of login] (reject1) {Rechazar solicitud};
\node[boxstep,right=2.7cm of elig] (reject2) {Deshacer transacción y rechazar};

\draw[edge] (start)--(login);
\draw[edge] (login)--node[right]{sí}(enc);
\draw[edge] (login.east)--node[above]{no}(reject1.west);
\draw[edge] (enc)--(db);
\draw[edge] (db)--(elig);
\draw[edge] (elig)--node[right]{sí}(save);
\draw[edge] (elig.east)--node[above]{no}(reject2.west);
\draw[edge] (save)--(end);

\end{tikzpicture}
}
\caption{Validación de la sesión, cifrado local, transacción atómica y emisión del recibo.}
\end{figure}

\subsection{Instrucciones de uso}
\textbf{Para MacOS}
\begin{lstlisting}
python3 -m venv .venv
source .venv/bin/activate
python -m pip install -r requirements.txt
python manage.py init
python launcher.py
# Abrir http://127.0.0.1:5000/login

# En otra terminal, al concluir:
python manage.py close
python manage.py tally
python manage.py verify
python -m pytest -q
\end{lstlisting}

\newpage
\textbf{Para Windows}
\begin{lstlisting}
#1. Construir el entorno virtual
py -m venv .venv

#2. Activarlo
.\.venv\Scripts\Activate.ps1

#3. Instalar las bibliotecas (Si es la primera vez que se ejecuta)
python -m pip install -r requirements.txt

#4. Inicializar la urna
python manage.py init

#5. Iniciar el servidor
python launcher.py

#6. Abrir en el navegador (preferiblemente Chrome) http://127.0.0.1:5000/login
bruh...
por definir
\end{lstlisting}
Los comandos administrativos se ejecutan \emph{localmente}; no se expone un botón web para cerrar o escrutar. 
